\section{The normalized factorization bridge}\label{sec:factorization-bridge}

The geometric construction admits a short affine realization before the
expanded polynomial is written down.  Let
\[
 L=aT+bS,
 \qquad
 Q=cT^2+dTS+eS^2.
\]
Then
\[
 LQ=(ac)T^3+(ad+bc)T^2S+(ae+bd)TS^2+(be)S^3
\]
and
\[
 \operatorname{Res}(L,Q)=a^2e-abd+b^2c.
\]

\begin{proposition}[coefficient--resultant coordinates]
The morphism
\[
 \Theta:\A^2\times\A^3\longrightarrow\A^4\times\A^1,
 \qquad
 (L,Q)\longmapsto\bigl(LQ,\operatorname{Res}(L,Q)\bigr)
\]
is etale on the nonzero-resultant locus.
\end{proposition}

\begin{proof}
Suppose that a tangent vector $(\dot L,\dot Q)$ is killed by the differential
of multiplication.  Then
\[
 \dot LQ+L\dot Q=0.
\]
When the resultant is nonzero, $L$ and $Q$ are coprime.  Hence $L$ divides
$\dot L$, and degree considerations give
\[
 (\dot L,\dot Q)=\lambda(L,-Q).
\]
This is the infinitesimal relative-scaling direction.  The resultant has
bidegree $(2,1)$, so
\[
 d\operatorname{Res}_{(L,Q)}(L,-Q)
 =(2-1)\operatorname{Res}(L,Q)\ne0.
\]
Thus $d\Theta$ has trivial kernel; source and target both have dimension five.
In the displayed coordinates the same statement is
\[
 \det D\Theta=-\operatorname{Res}(L,Q)^2.
\]
\end{proof}

Consider the normalized slice
\[
 X_{\mathrm{fac}}=
 \left\{(a,b,c,d,e):
 a^2e-abd+b^2c=1,\quad ad+bc=1
 \right\}.
\]
The multiplication map
\begin{equation}\label{eq:factorization-multiplication}
 \mu:X_{\mathrm{fac}}\longrightarrow\A^3,
 \qquad
 (a,b,c,d,e)\longmapsto(ac,ae+bd,be)
\end{equation}
is the base change of $\Theta$ to the slice on which the resultant and the
$T^2S$ coefficient are both one.  It is therefore etale, without a separate
calculation at $a=0$.

\begin{proposition}[polynomial slice coordinates]\label{prop:factorization-slice}
The morphism $\Phi:\A^3_{a,y,z}\to X_{\mathrm{fac}}$ defined by
\begin{align*}
 b&=1+ay,\\
 c&=1-\frac32ay+a^2z,\\
 d&=\frac12y-az+\frac32ay^2-a^2yz,\\
 e&=-2z+4y^2-4ayz+3ay^3-2a^2y^2z
\end{align*}
is an isomorphism.  Its inverse is
\begin{equation}\label{eq:factorization-inverse}
 y=2bd-ae,
 \qquad
 z=2d^2+ce+6bd^2+3bce-\frac92e.
\end{equation}
\end{proposition}

\begin{proof}
Substitution of the forward formulas gives
\[
 ad+bc=1,
 \qquad
 a^2e-abd+b^2c=1.
\]
Substituting them into \eqref{eq:factorization-inverse} returns $y,z$.
Conversely, substitution of \eqref{eq:factorization-inverse} into the forward
formulas returns $b,c,d,e$ modulo the two defining equations of
$X_{\mathrm{fac}}$.  All formulas are polynomial, including on the divisor
$a=0$.
\end{proof}

The residual torus action
\[
 \lambda\cdot(a,b,c,d,e)
 =(\lambda a,b,c,\lambda^{-1}d,\lambda^{-2}e)
\]
becomes
\[
 \lambda\cdot(a,y,z)=(\lambda a,\lambda^{-1}y,\lambda^{-2}z).
\]
Thus the weights $(1,-1,-2)$ are the relative factor-scaling weights on the
normalized slice.

Put
\[
 G=\mu\circ\Phi:\A^3_{a,y,z}\longrightarrow\A^3.
\]
Exact expansion gives
\[
 \det DG=-1.
\]
The polynomial $F$ of Section~\ref{sec:explicit} is obtained from $G$ by
linear source and target changes.  Namely, with
\[
 A(z_1,z_2,z_3)=\left(z_1,z_2,-\frac12z_3\right),
 \qquad
 B(u_1,u_2,u_3)=(u_3,2u_2,2u_1),
\]
one has
\begin{equation}\label{eq:factorization-to-foundational}
 F=B\circ G\circ A.
\end{equation}
Consequently
\[
 \det DF=(\det B)(\det DG)(\det A)=-2.
\]

This separates the three mechanisms.  Etaleness comes from coprimality and
the nonzero resultant weight; generic degree three comes from the three
choices of linear factor of a squarefree cubic; and the exceptional geometric
input is Proposition~\ref{prop:factorization-slice}, identifying this tangent
nonosculating slice with affine three-space.
